%  LaTeX support: latex@mdpi.com 
%  For support, please attach all files needed for compiling as well as the log file, and specify your operating system, LaTeX version, and LaTeX editor.

%=================================================================
\documentclass[data,article,submit,pdftex,moreauthors]{Definitions/mdpi} 

%=================================================================
% MDPI internal commands
\firstpage{1} 
\makeatletter 
\setcounter{page}{\@firstpage} 
\makeatother
\pubvolume{1}
\issuenum{1}
\articlenumber{0}
\pubyear{2022}
\copyrightyear{2022}
%\externaleditor{Academic Editor: Firstname Lastname}
\datereceived{} 
\dateaccepted{} 
\datepublished{} 
%\datecorrected{} % Corrected papers include a "Corrected: XXX" date in the original paper.
%\dateretracted{} % Corrected papers include a "Retracted: XXX" date in the original paper.
\hreflink{https://doi.org/} % If needed use \linebreak
\graphicspath{{figures/}} % path to folder with figures
%\doinum{}
%------------------------------------------------------------------
% The following line should be uncommented if the LaTeX file is uploaded to arXiv.org
%\pdfoutput=1

%=================================================================
% Add packages and commands here. The following packages are loaded in our class file: fontenc, inputenc, calc, indentfirst, fancyhdr, graphicx, epstopdf, lastpage, ifthen, lineno, float, amsmath, setspace, enumitem, mathpazo, booktabs, titlesec, etoolbox, tabto, xcolor, soul, multirow, microtype, tikz, totcount, changepage, attrib, upgreek, cleveref, amsthm, hyphenat, natbib, hyperref, footmisc, url, geometry, newfloat, caption

%=================================================================
%% Please use the following mathematics environments: Theorem, Lemma, Corollary, Proposition, Characterization, Property, Problem, Example, ExamplesandDefinitions, Hypothesis, Remark, Definition, Notation, Assumption
%% For proofs, please use the proof environment (the amsthm package is loaded by the MDPI class).

%=================================================================
% Full title of the paper (Capitalized)
\Title{Handling Dataset with Geophysical and Geological Variables on the Bolivian Andes by Scripts}

% MDPI internal command: Title for citation in the left column
\TitleCitation{Handling Dataset with Geophysical and Geological Variables on the Bolivian Andes by Scripts}

% Author Orchid ID: enter ID or remove command
\newcommand{\orcidauthorA}{0000-0002-5759-1089} % Add \orcidA{} behind the author's name
%\newcommand{\orcidauthorB}{0000-0000-0000-000X} % Add \orcidB{} behind the author's name

% Authors, for the paper (add full first names)
\Author{Polina Lemenkova $^{1}$\orcidA{}}
% \dagger,\ddagger
%\longauthorlist{yes}

% MDPI internal command: Authors, for metadata in PDF
\AuthorNames{Polina Lemenkova}

% MDPI internal command: Authors, for citation in the left column
\AuthorCitation{Lemenkova, P.}
% If this is a Chicago style journal: Lastname, Firstname, Firstname Lastname, and Firstname Lastname.

% Affiliations / Addresses (Add [1] after \address if there is only one affiliation.)
\address{%
$^{1}$ \quad Université Libre de Bruxelles, École polytechnique de Bruxelles (EPB, Brussels Faculty of Engineering), Laboratory of Image Synthesis and Analysis, Building L, Campus de Solbosch CP131/3, Avenue Franklin D. Roosevelt 50, B-1050, Brussels, Belgium; polina.lemenkova@ulb.be or pauline.lemenkova@gmail.com}

% Contact information of the corresponding author
\corres{Correspondence: polina.lemenkova@ulb.be; Tel.: +32471860459 (P.L.)}

% The commands \thirdnote{} till \eighthnote{} are available for further notes

%\simplesumm{} % Simple summary

%\conference{} % An extended version of a conference paper

% Abstract (Do not insert blank lines, i.e. \\) 
\abstract{This study presents processing georeferenced data for mapping geophysical and geological parameters of the Bolivian Andes, South America. The data included the following open sources: the IRIS seismic measurements (2015 to 2021), geophysical data from satellite-derived gravity grids based on CryoSat, topographic GEBCO data, geoid undulation data from EGM-2008, and geological georeferences vector layers from the USGS. The data were combined to form a dataset aimed at geological analyses of the Bolivian Andes. A highly seismically active region located in the South American Andes, this region requires a data integration for a detailed analysis of geophysical and geological setting. The techniques of data processing includes QGIS and GMT scripting toolset for quantitative and qualitative evaluation of the seismicity and geophysical setting in Bolivia. The result includes a series of seven geographic maps on Bolivian Andes. Based on the data analysis, the western region was identified as the most seismically endangered area in Bolivia with high risk of earthquake hazards in Cordillera Occidental, followed by the Altiplano and Cordillera Real. The earthquake magnitude here ranges from 1.8 to 7.6. The data analysis shows a tight correlation between the gravity and topography in the Andes. The utility of scripting cartographic techniques for geophysical and topographic data processing, combined with GIS spatial evaluation of the geological data, supported automated mapping that has applicability for risk assessment of the Bolivian Andes and geological hazard mapping in South America. The cartographic scripts used for processing data in GMT are available in the author's public GitHub repository in open access with provided link.}

% Keywords
\keyword{cartography; spatial data; GIS; georeferencing; mapping; Latin America; geodesy} 

% The fields PACS, MSC, and JEL may be left empty or commented out if not applicable
%\PACS{J0101}
%\MSC{}
%\JEL{}

%%%%%%%%%%%%%%%%%%%%%%%%%%%%%%%%%%%%%%%%%%
% Only for the journal Diversity
%\LSID{\url{http://}}

%%%%%%%%%%%%%%%%%%%%%%%%%%%%%%%%%%%%%%%%%%
% Only for the journal Applied Sciences:
%\featuredapplication{Authors are encouraged to provide a concise description of the specific application or a potential application of the work. This section is not mandatory.}
%%%%%%%%%%%%%%%%%%%%%%%%%%%%%%%%%%%%%%%%%%

%%%%%%%%%%%%%%%%%%%%%%%%%%%%%%%%%%%%%%%%%%
% Only for the journal Data:
%\dataset{DOI number or link to the deposited data set in cases where the data set is published or set to be published separately. If the data set is submitted and will be published as a supplement to this paper in the journal Data, this field will be filled by the editors of the journal. In this case, please make sure to submit the data set as a supplement when entering your manuscript into our manuscript editorial system.}

%\datasetlicense{license under which the data set is made available (CC0, CC-BY, CC-BY-SA, CC-BY-NC, etc.)}

%%%%%%%%%%%%%%%%%%%%%%%%%%%%%%%%%%%%%%%%%%
% Only for the journal Toxins
%\keycontribution{The breakthroughs or highlights of the manuscript. Authors can write one or two sentences to describe the most important part of the paper.}

%%%%%%%%%%%%%%%%%%%%%%%%%%%%%%%%%%%%%%%%%%
% Only for the journal Encyclopedia
%\encyclopediadef{Instead of the abstract}
%\entrylink{The Link to this entry published on the encyclopedia platform.}
%%%%%%%%%%%%%%%%%%%%%%%%%%%%%%%%%%%%%%%%%%
\begin{document}

%%%%%%%%%%%%%%%%%%%%%%%%%%%%%%%%%%%%%%%%%

\section{Introduction}

\subsection{Background}

For many problems in Earth sciences, cartographic datasets are the key sources for spatial analysis based on data visualisation. However, the most important issues in mapping workflow consists in finding the algorithms of data processing and solutions for processing multi-source and heterogeneous data. This paper proposes the integrating use of techniques of Geographic Information System (GIS) and Generic Mapping Tools (GMT) cartographic scripting toolset for processing geospatial vector and raster data.

Having the efficient scripting algorithms in cartographic data processing of geospatial datasets can bring significant speed improvements to mapping process \cite{CHEN2022105362,ORTOLANO201856, Lemenkova202203,OMRAN2016140,data5010008}. Examples of such applications include applying the modules of GMT for specific cartographic tasks of data processing and visualisation. Effective processing of geospatial data by scripts is especially actual for the large datasets with global coverage on Earth observation \cite{LEI2020202,XIA2018245,NATIVI20151,WAGEMANN2021104916}. 

The console-based data processing using scripts enables to increase the accuracy and speed of data processing in cartographic analysis and plotting. Therefore, using advanced geospatial software and tools is also of major importance in many GIS-related applications that includes geospatial data processing, such as machine learning applications in GIS \cite{1607434,ARABAMERI2021100848,ZUO2021105111,SUN201926}, data retrieval, exploration and processing \cite{BUCK2022121,Lemenkova202135,1369823}, spatial data formatting, handling and compression \cite{ZHU2022104166,data4030118,data6080081,DEMIRALTINTAS2021103743}, geographic data modelling and analysis \cite{MA2021103858,Lemenkova202205,TAMIMINIA2020152}. 

Processing geographic data is important with regard to a variety of the existing data formats. The two major geographic data formats, raster and vector types, are often used combined in a mapping workflow. However, the fundamental differences in the data structure entailed various means of handling and processing of these data.

Raster data and images are often derived from the remote sensing (RS) sources as satellite images. Examples of using raster geographic data in Earth related studies include satellite imagery. For instance, the Sentinel-1 and 2A, used in a variety of cases for mapping high-resolution spatio-temporal dynamics of land use and land cover \cite{MERCIER2021107408,BALENZANO2021107345}. The wide applications of the Landsat TM/ETM+ by ESA for geological and land cover mapping are caused by the data availability and global coverage \cite{ALLENBY1987137,LOMBARDO,SAMARA}. The aerial Google Earth or the commercial geospatial data providers operating of large-scale aerial data present large-scale data for detailed mapping. More examples of raster data sources include hyperspectral imagery, InSAR largely used in geology related studies \cite{TALIB2022112793,LI2022114}, or LiDAR used for geomorphic computations or vegetation canopy height \cite{NEUGIRG2016269,NAKAO2022119953}, Important geospatial data sources for topographic mapping are given by the SRTM DEM. 

Using raster data as a source of spectral information is often a case for land cover mapping and landscape analysis \cite{SZABO2019104187,RUJOIUMARE2016244,LI2021112670}. Scanned raster images in bitmap formats (jpg, bmp, tiff) provide yet another source of geospatial information aimed at retrospective data analysis used for geospatial reconstructions \cite{KRAL2021107382,OSTAFIN2022107709}. Used as georeferenced raster layers in GIS, archived historical maps can be utilized as data sources for estimating changes level to evaluate landscape dynamics and serve as auxiliary source for geologic analysis and mineral resource management. Technically, processing raster  data is based on the analysis of pixels on the image and, if necessary, steps of vectorisation \cite{8937640,6723417,8810063}. In the advanced cases, the RS data can be processed using the machine learning techniques. 

Vector data, on the other hand, are often used for representing geographical information on the object-based level \cite{LI20121331,KOCH200623}. The information containing as ontologies of vectors data is included in the attribute tables of the associated GIS tables and/or metadata of the layers \cite{LI20121344}.
Vector data sources include, \emph{e.g.}, the vectorised objects as georeferenced layers in GIS formats (\emph{e.g.}, the well known .shp format of the ArcGIS/QGIS). The integration and complementary use of vector to raster data can be illustrated by the cases of cadastral cartography and LiDAR altimetric data providing information for the 3D vector-based model of the architecture structures \cite{VIANAFONS2020110258}. A special feature of vector data consists in object-based structure which entails points, lines or polygons formed by a variety of curves or points with coordinates in a spatial reference format. Thus, vector objects indicate boundaries or depict the contours of geographic objects that vary in geometry, form, size, shape, curvature, connectivity and spatial resolution, depending on map scale \cite{GULER2021107755}.

Although existing GIS for geospatial data processing are widely available, directly using their functionality does not offer an intuitive, direct control on data handling as in case of using scripts. In fact, various GIS software often provide the designed solutions for mapping workflow and processing data. In contrast, script-based technique provides more control on data processing. This article present an integrative data-driven cartographic synthesis technique that enables processing vector geologic and raster geophysical data, while exploring all the advantages of GIS and GMT applied in one framework.

This study is focused on processing the spatial data aimed for cartographic visualisation using geological, topographic and geophysical datasets on Bolivia (Figure~\ref{fig1}). Integration of the vector and raster data in a structured dataset by means of various technical tools supports the increase of geographical and expert knowledge. Creating such a project with a special focus on Bolivian Andes is built upon the data captured from the available open sources. The framework of this study is designed to visualise data from various sources, ontologies and formats as a series of maps aimed at monitoring natural resources in a selected region of the geographically complex region of South America.

%\newpage
\subsection{Stuy Area}
Identifying correlations between geological structure of the land, its topography and geophysical setting and earthquake magnitude, depth and frequency gives better insights into our understanding of the geological hazards and risk assessment in the mountainous regions. The study object is the Bolivian segment of Andes Mountains. The Andes Mountains are a unique geologic, geomorphological and tectonic phenomena built upon the thickest crust on Earth. Andes are the longest continental mountain range that stretches along the western margins of the South America and crosses several countries including Bolivia (Figure~\ref{fig1}).

\begin{figure}[H]\centering
	\includegraphics[width=12.0 cm]{Fig_01}
	\caption{Topographic map of the study area. Visualization: GMT. Source: author.\label{fig1}}
\end{figure}   
%\unskip

Andes are the key geological object of Bolivia directly affected by its tectonic history and Earth’s history. The geological evolution of the Bolivian Andes has underwent a long multi-staged period of formation which has been reported and described in details in previous studies \cite{Andersonetal,Arguedas} and in some earlier fundamental books on the South American geology \cite{Forbes,Suarez,Ahfeld}. Bolivia occupies the Central Andes segment with the Altiplano, the world's 2nd highest mountain plateau after the Tibet. The tectonics of the Andes largely control the geomorphology of the Bolivian mountainous lakes (shape, orientation and geometry), primarily by a nearly orthogonal fracture pattern from the underlying granitic basement (Brazilian Shield), and secondary, by the tectonic movements: ancient global fracture pattern, principal compressive stress, downwarping and uplifting of the Beni Basin during the Cenozoic Andean orogeny \cite{Allenby1988}. 

The specific features of tectonics and geology of Central Andes and seismicity have been described by many researchers \cite{Ramos,Fernandez,Funning} and have drawn the attention of Earth scientists in various disciplines, e.g. hydrogeologists \cite{Lima2021a,Lima2021b}, geophysicists \cite{Aviles,Borsa}, seismologists \cite{Dorbath,Dahm}. The historic evolution of the Bolivian Andes underwent several stages during the early history of the Earth. Among the notable geologic events of Bolivia the San Ignacio Orogeny which reshaped the terrane through deformations accompanied by voluminous tectonic granite intrusion and deformation of the sediments and intrusions \cite{Boger}. 

\begin{figure}[H]\centering
	\includegraphics[width=14.0 cm]{Fig_02}
	\caption{Geologic map of Bolivia and surrounding countries. Visualization: QGIS. Source: author.\label{fig2}}
\end{figure}   

The dominating units in central Bolivia (Figure~\ref{fig2}) are presented by the Cambrian-Ordovician (CmO) outcrops (dark magenta) which largely contrasts with a very complex disposition of several layers in the west of the country in Central Andes region: Ordivician-Silurian (OS), Devonian (D), Silurian (S), Tertiary (T), and Creatatios-Tertiary volcanics (Cv). The eastern part of the country (the Amazon forests) are occupied by the Precambrian outcrops (pC) which well corresponds with the extend of the Brazilian Shield (Figure~\ref{fig2}). 

General trend of the geological structures of the Bolivian orogeny is oriented at the NW direction off Amazonia following the extent of the Central Andes. The most important geological provinces of Bolivia (Figure~\ref{fig2}) include Madre dos Dios on the north, mostly covered by theTertiary rocks (T), the Brazilian Shield on the east, the Oran-Olmedo Basin on the south (covered by the Cambrian-Ordovician outcrops), the Beni Basin in the central-northern segment of the country, and the Andean Province on the west (covering the Central Andes) with intrusion of the Altiplano Basin (Figure~\ref{fig2}). 

Andes act as a source of precious mineral and geological resources in Bolivia due to the intensive tectonic processes and orogenesis in its basins and geologicial provinces (Figure~\ref{fig3}). The formation of metals in Andes includes complex petrographic and geochemical processes that results in the evolution of magmas and lavas that underwent fractionation and assimilation in the Earth’s crust \cite{Redwood}. As a result, mining industry Bolivia has a dominant position in the country’s economy since 16th century, which resulted in its leading position in the world in terms of reserves of tin, zinc, silver, copper, gold, tungsten, lead, antimony  \cite{Iqbal,DiazCuellar,Guedron,Pavilonis}. The Bolivian Salar de Uyuni salt flat is the world's largest lithium deposit \cite{Narins,Sanchez}. 

\begin{figure}[H]\centering
	\includegraphics[width=14.0 cm]{Fig_03}
	\caption{Geologic provinces in Bolivia and surroundings. Visualization: QGIS. Source: author.\label{fig3}}
\end{figure} 

Regional geomorphology of Andes affects a wide variety of interconnected setting of Bolivia: hydrology, environment, soils, ecology, climate and landscapes. For instance, the tectonic faults, lineaments and structure of the underlying rocks cause variations in the geomorphic and hydrological fluvial structure of Bolivia. This was recorded within the drainage network of the Andean foreland which are preserved in the paleochannels of the Bolivian Amazon that vary in sinuosity, widths and pattern of the Amazon rivers \cite{Plotzki}. Landscapes of Bolivia vary according to their location within Andes which serve as a natural barrier to the air masses controlling climate, causing fluctuations in temperature and humidity and water balance in Bolivia. 

Various regions of the Bolivian Andes give place to the National Parks and protected nature conservation zones \cite{Baumannetal2018,Thomasetal,Burbridge}. Besides mountainous regions, the forests of the Bolivian Amazon provide ecological habitat for the endemic species \cite{Cadima} which demonstrates the environmental significance of this region. Besides climate change and natural factors affecting the landscapes, anthropogenic activities contributed to the landscape changes in Bolivia. Thus, according to the study in Iroco and Cueva Bautista regions in the highlands of Bolivia, human settlements in the high-altitude ecosystems (>3800 m asl) occurred, at least, by 13,000 years BP \cite{Capriles}.  

At the same time, the seismicity of Andes is high with repetitive earthquakes of various focal depth, rupture size and magnitude \cite{Rivadeneyra,Tono}. In turn, earthquakes are closely associated with tectonic activity and zones of plate subduction. Magma evolution and assimilation at the subduction zones of Andes significantly affected the geologic setting in South America \cite{Gonzalezetal2020,Lemenkova2019a} 2020e). Earthquakes usually have negative consequences and involve such hazardous events as landslides, rockfalls, ground shaking that results in damaged and destroyed buildings, ruined bridges and other constructions, affected or lost human lives, general losses in economics and financial instability. This resulted in intensive studies for Central Andes and geological investigations of Bolivia.

%%%%%%%%%%%%%%%%%%%%%%%%%%%%%%%%%%%%%%%%%%
\section{Materials and Methods}

Methods of GIS mapping and geoinformatics passed several epochs from its onset in 1970s, including rapid development of computer technologies, progress in spatial data standards and formats, remote sensing and processing of Earth observation and geological data \cite{Allenby1987,Lindh2021,Lemenkova202127}, approaches of semi-handmade to machine-based cartographic data visualization \cite{SchenkeLemenkova}. Nowadays, the contemporary methods of spatial data processing include high degree of computing and scripting with recent application of machine learning, scripting and programming in geosciences \cite{Lietal,Lemenkov2021a,Zhou,Lemenkov2021b}. 

Current study was conducted to identify the distribution and magnitude of earthquakes in the Bolivian Andes in context of its geologic, geophysical and topographic situation. The study was performed using the technical integration of QGIS and GMT scripting tools that were applied to assess the geophysical setting in Bolivia against topographic elevation data, geologic provinces and units with plotted data of seismic events obtained from the IRIS database. 	 

The topographic map (Figure~\ref{fig1} and Figure~\ref{fig7}) was mapped as a background and against an IRIS seismic data from General Bathymetric Chart of the Oceans (GEBCO) gridded bathymetry and topography dataset \cite{Schenke2016}. GEBCO presents consistent and accurate digital elevation data covering the Earth and largely used in Earth science and reliable data source. The unprecedentedly high resolution (15 arc second) of GEBCO grid enables to use it as a fine cartographic grid both for topographic mapping and as a background for overlaying thematic maps \cite{Lemenkova2020b}, 2020c). 

This study uses the EGM-2008 \cite{Pavlis} digital geoid model to show that there are variations in gravity over the territory of Bolivia that in general well correlate with the topographic map (compare Figure~\ref{fig1} and Figure~\ref{fig4}) with a general increase of gravity values in the south-west direction that coincides with the extend of Central Andes. The data on the gravity (Figure~\ref{fig5} and Figure~\ref{fig6}) were obtained as remote sensing grids from the repositories of Scripps Institution of Oceanography \cite{Sandwelletal2014}. 

The IRIS database of the world seismicity applied for Bolivia for the period of 2015–2021 was used to map the earthquake events recorded in central Andes for the six-year time span (Figure~\ref{fig7}). Several columns of the initial table in the CSV format presented seismic parameters of the earthquakes (depth and magnitude), coordinates and time of the event (date-month-years), short description and location. The studied table included 3,000 records with seismic events in Bolivia. The table was converted from CSV to the NGDC format and used for GMT plotting (Figure \ref{fig7}). The study employed existing GIS methods \cite{Hubbard,Lemenkova202135,Suetova,Lemenkova202125,Klauco2013,Klauco2017} and presented two maps (Figures 2 and 3) made using the QGIS software \cite{QGIS}.

Besides GIS, the study applied scripting Generic Mapping Tools (GMT) developed and continuously maintained by \cite{Wesseletal2019} for mapping in Figure~\ref{fig1} and Figures~\ref{fig4}--~\ref{fig7}. In contrast to the traditional GIS, the GMT refers to a console-based scripting framework of coding principle for automated plotting of the geospatial data. The GMT cartographic approach is an effective automated way of mapping and data visualization which is based on shell scripting \cite{Gauger,Lemenkova2020a,Lemenkova2020d}, and includes a variety of data processing steps such as conversion, reshaping and reformatting, data modelling, spatial analysis and visualization. The process of the GMT mapping was adapted based on the research performed in previous works \cite{Lemenkova202140}, 2021a).

%%%%%%%%%%%%%%%%%%%%%%%%%%%%%%%%%%%%%%%%%%
\section{Results}

In present study the geophysical settings and geological units of Bolivia were analysed for the comparison with seismicity and topographic elevations presenting the geomorphological variability and formed unique landscapes of the country. The topography of Bolivia (Figure~\ref{fig1}) shows regional variations that well reflect Quaternary geologic history that includes the periods of glaciations of higher mountains in the Bolivian Andes that were glaciated in the Pleistocene, e.g. Cordillera Real \cite{Heine}. Glaciation in Bolivian Andes experienced several stages of periods in its early geologic history that include Tertiary and Early Pleistocene glaciations (> 790 ka), Middle Pleistocene glaciations (790–132 ka); Late Pleistocene glaciations (132–11.7 ka); Holocene glaciations (11.7 ka–to present) according to \cite{Frenierre}. The results of the glaciations remained as bedded slope deposits on moraine ridges in the Bolivian Andes. 

\begin{figure}[H]\centering
	\includegraphics[width=12.0 cm]{Fig_04}
	\caption{Geoid model of Bolivia. Visualization: GMT. Source: author.\label{fig4}}
\end{figure}  

Late Pleistocene glaciations also formed megafans of Gran Chaco, a huge plain in eastern Bolivia and the main biogeographic biome of South America. Also contrasting glaciation formed large piedmont fans of Andes that were developed during colder and seasonal conditions and dry–wet intense seasons and enhanced rainfalls by orographic effects of the Andes \cite{Latrubesse2009,Latrubesse2012}. As a result, complex interrelation of the vegetation, climate and soils can be illustrated by the reduced forest cover as result of increased aridity in selected landscapes fo Chaco, soil formation and intensive sedimentation in the denser forest cover in the Bolivian Chaco \cite{May2008}. Besides glaciation, other external factors that strongly affected and shaped the topography of Bolivia include fluvial and aeolian interactions, climate and biogeographic effects. These mostly formed the Gran Chaco and regions of Amazonia through subtropical semi-deciduous vegetation that spreads on the Andes piedmonts together with soil and climate setting.

The geomorphological patterns initially formed by the technical and geological processes are mirrored in fluvial belts, lacustrine patterns and riverine meanders of Bolivia and create landscapes of the country \cite{Dumont}. Landscapes in eastern Bolivia located along the southern margin of the Amazon basin are affected by climate notable by strong hydrological and climatic seasonality, e.g. high wind speeds and aeolian processes. This affects the geomorphology and environment of eastern Bolivia forming special aeolian landforms: linear sand streaks and ridges, parabolic and elongate morphological structures, dunes, sand sheets, mega-ripples. hence, the climate and aeolian processes (wind, sand sediments) control the environment, vegetation and soils \cite{May2013}. 

\begin{figure}[H]\centering
	\includegraphics[width=12.0 cm]{Fig_05}
	\caption{Free-air gravity map of Bolivia. Visualization: GMT. Source: author.\label{fig5}}
\end{figure}   

The geologic maps (Figure~\ref{fig2} and Figure~\ref{fig3}) reflect the current state of the processes of solid Earth physics in the Andes. Thus, the metamorphic rocks and outcrops of various geologic epochs experienced solid-state mineral transformations due to changes in pressure and temperature in crust and mantle of the Earth, and record geological periods, heating, seismic processes and cooling that reflect the tectonic environment they were formed \cite{Holder}. As a result, complex geological history of Bolivian Andes reflects its regional geology and petrographic setting, as well as the origin, distribution and geochemistry of minerals \cite{Tapia}. 

In general, regional undulations of geoid (Figure~\ref{fig4}) well correlate with the topography of the country. Thus, the values of geoid clearly increased in westward direction exceeding 34 m, which well correlates with the geomorphology of the Central Andes. Here the isolines increase towards the west of Bolivia. In contrast, negative values of geoid (from 10 to -10) are visible in the NE region of the study area and mostly cover the Amazon region of Brazil. The Beni Basin of Bolivia (Figure~\ref{fig3}) demonstrates slight decrease in data (from 20 to 16 m) while the central regions of the Madre dos Dios Basin shows local increase in geoid values (up to 32 m in the surroundings of the Madidi National Park).

\begin{figure}[H]\centering
	\includegraphics[width=12.0 cm]{Fig_06}
	\caption{Vertical gravity gradient in Bolivia. Visualization: GMT. Source: author.\label{fig6}}
\end{figure}   

A distinct correlation shows the comparison of geological map in Figure~\ref{fig2} with geoid model in Figure~\ref{fig4}.  The distribution of the geologic folds westward off the Cambrian-Ordovician outcrops well repeats the dense and steep increase in isolines of the geoid detecting the increase of gravity that well correlates with the mountain areas of the Andes (red colors areas in Figure~\ref{fig4} shows values of geoid over 35 m). The correlation in changes of the free-air gravity (Figure~\ref{fig5}) and vertical gravity gradient (Figure~\ref{fig6}) with topographic map (Figure~\ref{fig1}) and geologic map (Figure~\ref{fig2}) are found by consecutive comparing these maps, further from Andes to Amazonia area and the Gran Chaco region. In the inner western region of the Andes (Cordillera Occidental) a series of volcanoes are detected and mapped (Figure~\ref{fig7}) that show correspondence with higher magnitude of earthquakes (green dots indicating M 3.4 to 4.0 and 4.1 to 4.7). 

The investigation of the values of the free-air gravity (Figure~\ref{fig5}) and vertical gravity (Figure~\ref{fig6}) demonstrates very good correlation with local lacustrine and riverine depressions of the mountainous lakes in Bolivia (Laguna Rogagua, Laguna San Luis), rivers of Amazonia, regional depression of Gran Chaco on the border with Paraguay, and the piedmonts of the Andes (dark blue colors in Figures~\ref{fig5} and ~\ref{fig6}). In contrast, the highest elevation peaks of the Bolivian Andes (Cordillera Real, Cordillera Occidental and the Altiplano) well correlate with the highest values in geophysical data. The correlation of the gravity data with the geoid (Figure~\ref{fig4}) well exhibit the lowest geoid values in the regions of the lower gravity in Figures~\ref{fig5} and ~\ref{fig6} and the saddle-like local depression in the region of the city La Santísima Trinidad (central Bolivia) showing local topographic depression. 

\begin{figure}[H]\centering
	\includegraphics[width=12.0 cm]{Fig_07}
	\caption{Seismicity in Bolivia and the Central Andes. Visualisation: GMT. Source: author.\label{fig7}}
\end{figure}   

The overall distribution of the earthquakes in Bolivia (Figure~\ref{fig7}) indicated a good correlation with geologic extent of the Cordillera Occidental, followed by the Altiplano and the Cordillera Real. In total, 3,000 records have been assessed using IRIS database for years 2015-2021. In general, the higher seismicity level of Bolivia is clearly visible in the westernmost region of the country. The earthquake with the least magnitude (1.8) is detected 17 km ESE of Iquique, Chile at depth 20 km. In contrast, the most significant, very deep earthquake events are located 211 km S of Tarauaca, Brazil (620.6 km, magnitude 7.6), and at 173 km WNW of Iberia, Peru with the depth of 606.2 km.

The most significant earthquake in situated in Bolivia is located 11 km NNE of Carandayti with depth of 559 km and magnitude of 6.8. Also, some significant earthquakes within the strongest magnitude (above 7) are located in Peru: 1) 140 km WNW of Iberia, Peru which took place in 2018 with magnitude of 7.1 at depth of 630 km; 2) 38 km SSW of Acari, Peru, also dated of 2018, with magnitude of 7.1 located at depth of 39 km (shallow earthquake). Thus, one can conclude that in Bolivia, the most significant earthquake during the last 6 years (2015-2021) did not exceed magnitude of 7.0, which indicates a lower seismic activity in Bolivia and less risk to population compared to Peru. 

\section{Discussion}

Using scripting techniques of GMT has many applications in geosciences. For example, study geologic setting of the Earth, climate modelling, exploring the atmosphere-ocean-solid earth interactions by advanced mapping, and many more. Rapid processing of maps by GMT enables to produce a series of maps, which is a crucial source of geo-information for such tasks. To fully utilise the potential of GMT, the abilities and advantages of modern scripting approaches should be used. In this research I demonstrated the functionality of GMT and QGIS for thematic mapping of Andes region from the geological and geophysical datasets. Currently geo-information data are available online in many sources which increase the need for rapid and robust processing of these data, based on availability. 

To do this, I introduced a data-driven approach of GMT and QGIS for thematic geological-geophysical mapping of Andes region, Bolivia. Using the functionality of this program this paper presented a workflow where I showed that it performs well on cartographic data processing without manual tedious workflow usually applied in cartography. I extracted the geological features of regional objects in Andes, mapped them as layers using QGIS and derived the information containing geological structure of the country: geological units including outcrops of various geologic time: Precambrian-Devonian, Paleozoic, Carboniferous, Cambrian, Ordovician, Permian, Tertarym Jurassic, Mesozoic, Cenozoic, Silurian, Triassic and Quaternary as the major units. The algorithms of GMT processed well relatively large files (each raster file mapped for visualization of geoid, gravity and seismicity has a size of several hundreds of MB) and converted these data into GMT native format for processing. Some complex regions of seismicity with overlapped sources of earthquakes were visualised using various colours and sizes of legend units for a better interpretation. Graphical representations of maps are illustrated in relevant sections. 

The process of mapping multi-source geo-information data is based on processing and visualising relevant thematic layers showing geological and geophysical features of the target objects (free-air gravity anomalies, vertical gravity gradients, earthquakes and other concepts). The major principle of GMT is creating the decision on visualising maps through a lines of code based on the GMT syntax language. For instance, I set up experimentally the projection and selected colour palettes for each corresponding phenomena to better show the variations in levels and to indicate the correlations between the topography and geophysics of the Andes region, which was used by the use of QGIS and GMT to identify the relationships between the geological and geophysical phenomena of the Earth through the prepared maps. Based on the setup parameters of GMT, the machine mapped the series of maps using scripts run from the console by a step-by-step cartographic process applied for each from the series of the presented maps.

This workflow arrived at data processing and visualisation on the presented series of maps (geologic, geopjhysical and topographic maps of Bolivia). However, certain steps in cartographic workflow required manual adjusting and improvements in QGIS compared to GMT, which demonstrated a higher level of automation in geo-information data processing, as shown in this work. The most geologically complicated regions (folds in the Bolivian Andes) were mapped using enlarged zoom, to gain more insights in the visual characteristics of the small areas (see the Devonian and Ordovician outcrops in Figure~\ref{fig2}) and to discriminate the strata from the undifferentiated areas. To do so, I zoomed selected enlarged areas on the map in the QGIS environment and controlled visualisation manually using a zoom function of QGIS. Specifically, I identified densely interspersed outcrops of Cretaceous-Tertiary volcanics (Cv), Precambrian-Devonian (AD) and Paleozoic (PZ) units densely interspersed in the Bolivian Andes,  which demonstrated a certain correlation with gravity anomaly and topographic heights on the same territory. 

This study presented an integrated application of the GMT scripting toolset and QGIS software for mapping geological and geophysical setting of Bolivia in a semi-automatic regime for GMT and manual workflow in QGIS. I noted that while using a reliable software tool is a major task of any cartographic research, there are still more algorithms that need to be automated and adjusted in GMT for a more automated cartographic workflow and processing of geo-information. Processing geo-information data is an essential part of geophysical and geological research having its specifics compared to the topographic or environmental mapping: interpretation of geophysical anomalous fields, visualising seismic events as earthquakes, recognition of geoid undulations via isoline directions and interpretation, smooth data conversion from seismic IRIS catalogue into GMT, etc. The goal of this research was to present here a practical application in geo-information data processing by integrated use of GMT and QGIS using dataset on geological, topographic and geophysical datasets.

\section{Conclusion}

Cartography and the entire system of Earth sciences experienced dramatic and positive changes when the programming and scripting methods were introduced in spatial data processing and thus helped in getting new insights into the correlations between the Earth phenomena and previously unknown processes through complex and integrated data visualization, statistical plotting and modelling of Earth observation data \cite{Bar,Lemenkova2019b,LOMBARDO,Lemenkova2020e}. A variety of GIS applications are presented in studies on spatial data visualization, processing and mapping. For instance, evaluating landscape metrics was applied in the research focused on landscape ecology and environmental issues in landscape analysis by GIS and remote sensing \cite{Millington,Lemenkova202137}. 	

Geomorphological GIS-based mapping is used for integrated thematic maps in agricultural, hydrological and geo-archaeological applications or in risk assessment studies \cite{Vella,Lemenkova202120}. Geological GIS-based mapping largely supported by visualization of the provinces and units \cite{Nedel} or statistical graphs and plotting for hydrocarbon resource analysis in geological context \cite{Jemio}. Other approaches include geophysical graphical plotting for analysis of magnetism \cite{Collinson,Ocola} or gravimetric geoid modelling \cite{Corchete}. GIS applications in mapping Bolivia are presented in studies of ecosystem services combining quantitative and qualitative data and spatial GIS analysis to assess environmental changes \cite{Hinojosa}. 

In contrast to all these and others thematic applications of GIS in South American studies with focus on Bolivia, geomorphology and environment, using scripting techniques in cartography presented a variety of automated tasks and opportunities of machine learning to cartography. Specifically, scripting approach causes advantages over the traditional GIS in geologic mapping and geophysical visualization due to the increased speed of mapping, accurate plot visualization and machine-based cartographic routine that minimizes the human caused errors and misprints. These include such improvements as automated mapping, precise machine-based data processing, significantly increased speed of spatial data handling, automated re-projecting of geodetic data, big data processing (gigabytes of volumes). Hence, using scripting enables to meet high demands of contemporary cartography, e.g. spatial computing, visualization of multi-format data from geological, geophysical, topographic, geodynamic and seismic sources for integrated geospatial data analysis.


%%%%%%%%%%%%%%%%%%%%%%%%%%%%%%%%%%%%%%%%%%
\vspace{6pt} 

\dataavailability{Not applicable} 

\acknowledgments{The author thanks the reviewers and editor for the review and editing of this manuscript.}

\conflictsofinterest{The author declares no conflict of interest.} 

\abbreviations{Abbreviations}{
The following abbreviations are used in this manuscript:\\

\noindent 
\begin{tabular}{@{}ll}
GIS & Geographic Information System \\
GMT & Generic Mapping Tools \\
QGIS & Quantum GIS \\
GEBCO & The General Bathymetric Chart of the Oceans \\
EGM & Earth Gravitational Models \\ 
IRIS & Incorporated Research Institutions for Seismology \\
USGS & United States Geological Survey \\
CSV & comma-separated values \\
NGDC & National Geophysical Data Center
\end{tabular}}

%%%%%%%%%%%%%%%%%%%%%%%%%%%%%%%%%%%%%%%%%%


%%%%%%%%%%%%%%%%%%%%%%%%%%%%%%%%%%%%%%%%%%
\begin{adjustwidth}{-\extralength}{0cm}
%\printendnotes[custom] % Un-comment to print a list of endnotes

\reftitle{References}

% Please provide either the correct journal abbreviation (e.g. according to the “List of Title Word Abbreviations” http://www.issn.org/services/online-services/access-to-the-ltwa/) or the full name of the journal.
% Citations and References in Supplementary files are permitted provided that they also appear in the reference list here. 

%=====================================
% References, variant A: external bibliography
%=====================================
\bibliography{Andes.bib}

%%%%%%%%%%%%%%%%%%%%%%%%%%%%%%%%%%%%%%%%%%
\end{adjustwidth}
\end{document}

